% ============================================================
% DRAFT NOTE (abstract)
% Contains: The deployed-system framing, the current-state hazard target, the
% training corpus and transfer setup, and the headline result.
% Written now: Problem framing and objective under the 2026-09-01 pivot to
% current-state multi-horizon divergence forecasting.
% Pending: All empirical numbers and the final outcome statement.
% Sources: pivot discussion (current-state hazard objective); prior HERALD
% draft (voice); Jocana/herald-logits dataset card
% ============================================================
\begin{abstract}
As language models are deployed in longer-context and lower-latency settings,
inference-time approximations increasingly determine not only cost but
reliability. KV-cache compression makes long-context inference cheaper, but
today its damage is evaluated after generation, by averaging task scores over
completed outputs. This misses the question a deployed system actually needs
to answer: while decoding, is compression damaging this continuation right
now, and how much damage is coming in the next few tokens? We introduce
HERALD, an online forecaster of compression-induced damage built from the
model's own logits. At every decoding step, whether or not the cache is
currently compressed, HERALD predicts the divergence the compressed
generation will accumulate against the uncompressed model over the next $k$
tokens, evaluated on the same realized prefix so that ordinary autoregressive
drift is factored out. The forecaster is a causal multi-horizon model over
zero-cost statistics of the next-token distribution, trained on a public
corpus of per-token compressed-generation traces and evaluated against strong
temporal baselines, with grouped, stratified evaluation by task, compressor,
ratio, horizon, and position. We validate that predicted distributional
damage tracks realized task-quality failures, and we test cross-model
transfer from the training corpus to a second model family.
\todo{Insert the final held-out evidence: per-horizon correlations, baseline
comparisons, quality-failure validation, and the transfer result.}
\end{abstract}
